\chapter*{Nota de método}
\addcontentsline{toc}{chapter}{Nota de método}

Este documento declara sus convenciones antes del primer cuadro. El género al que
pertenece no acostumbra hacerlo, y esa es la primera diferencia deliberada.

\section*{Objeto: aprobado contra aprobado}

Las cifras de gasto son del \textbf{Presupuesto de Egresos aprobado} de 2025
contra el aprobado de 2024, tomadas de los analíticos presupuestarios. Las de
ingreso son de la \textbf{Ley de Ingresos aprobada} de cada año. Donde los
Criterios Generales de Política Económica o el decreto son la única fuente, se usa
el proyecto y se dice en el lugar.

Esto se aparta de la práctica del género, que evalúa el proyecto porque publica
setenta y dos horas después de la entrega. Nosotros escribimos con un año de
distancia y no tenemos esa restricción; a cambio, perdemos la oportunidad de decir
algo antes de que la Cámara vote. Donde el cambio legislativo importa, se reporta
en una columna propia.

\section*{Deflactor y añada}

El deflactor es el del PIB que declara el CGPE 2025 para ese año: \textbf{4.3\%},
es decir un factor de 1.043. No se deriva de agregados publicados redondeados.
Toda cifra real está en \textbf{pesos de 2025}. Las razones a PIB usan el PIB
nominal de cada año tal como lo publica el CGPE 2025: 33,927.7 miles de millones
para 2024 estimado y 36,166.4 para 2025.

\section*{Las cinco declaraciones}

Toda comparación de este documento declara cinco cosas: si es nominal o real;
contra aprobado o contra cierre; contra el año previo o contra el PIB; en pesos de
qué año; y con qué PIB de qué añada. Cuando alguna no se puede declarar, la
comparación no se hace.

\section*{Una regla que este documento estrena}

\textbf{Toda diferencia expresada en millones de pesos dice si es nominal o real.}
Los cuadros traen las dos columnas. La razón es concreta: en el ejercicio 2024, un
documento del género presentó ocho de ocho variaciones en millones de pesos en
términos reales sin decirlo, y en un caso eso produjo una afirmación falsa sobre
un monto aprobado por la Cámara de Diputados. La convención de trabajar en pesos
constantes es defendible; no declararla no lo es.

\section*{Perímetros}

Cada capítulo declara su perímetro antes de su primer cuadro y no lo cambia dentro
del capítulo. Los perímetros que este documento usa y que no son evidentes:

\begin{itemize}
\item \textbf{Pensiones y jubilaciones}, clasificación económica: tipo de gasto 4
de las entidades de control directo, más el tipo de gasto 4 del Gobierno Federal
excluyendo la partida 45203, que es la transferencia del Gobierno Federal a las
entidades y se contaría dos veces. El perímetro no se eligió: se adjudicó. El
Anexo 3 del decreto publica los gastos obligatorios con y sin pensiones, y su
diferencia reproduce esta construcción al millón de pesos en los dos años.
\item \textbf{Salud}: función 2.3, Gobierno Federal más entidades, en bruto. El
perímetro canónico de la casa se ancla al cuadro de la exposición de motivos del
proyecto de presupuesto, documento que no existe para los ejercicios 2022 a 2026
porque está caído en el servidor de la Secretaría. La sustitución de fuente cambia
el objeto: los analíticos son brutos y aquel cuadro era neto. Se declara en el
capítulo.
\item \textbf{Educación}: función 2.5, Gobierno Federal, en bruto. Las entidades
de control directo tienen cero en esta función.
\item \textbf{Seguridad}: quince subfunciones de cinco funciones de la finalidad
Gobierno, enumeradas en el capítulo para que el lector pueda rehacerlo.
\end{itemize}

\section*{Proveniencia de las cifras del CGPE}

\textbf{Los Criterios Generales de Política Económica 2025 se publicaron sin capa
de texto.} Noventa de sus noventa y una páginas son imagen; la única con texto
contiene las fórmulas del PIB potencial. Se comprobó que el portal de la
Secretaría no sirve otra edición.

En consecuencia, toda cifra de los Criterios en este documento se leyó de la
página renderizada, y \textbf{no se validó releyéndola sino por identidad
contable}: cincuenta pruebas de cierre y de coherencia entre anexos, todas las
cuales pasan. La lista está en el anexo metodológico. Es el mejor control
disponible y no equivale a una restitución de texto; el lector debe saberlo.

\section*{Tier de cada cifra}

Toda cifra lleva su procedencia registrada: \emph{oficial primaria} para los
documentos del paquete y los analíticos, \emph{autoral ITED} para las
construcciones propias, y \emph{derivada} para lo tomado de terceros, siempre con
documento y página. El registro completo vive en el archivo de fuentes que
acompaña a los datos.

\section*{Lo que este documento no hace}

No presenta el marco de manual de sostenibilidad, primario más tasa menos
crecimiento. Se retiró por sesgo sistemático: el paquete no publica un balance
primario de los requerimientos financieros, y el del balance económico sesga al
optimismo por una o dos décimas. El diagnóstico interno queda en la bitácora del
proyecto.

\clearpage
